% title: 2026 春复变函数 (H) 期中考试
% author: 沈维孝
% summary: 从香蕉空间旧站“试卷：复变函数”迁移的完整试卷题目。
% source: https://bananaspace.org/legacy/wiki/用户:Solution/_试卷:_复变函数
\section*{2026 春复变函数 (H) 期中试卷（沈维孝）}
\begin{enumerate}
\item[一]（20分）设\(C : |z| = 4\)为逆时针定向圆周，计算
\[
    \int_C \frac{z^3 + 3z + 1}{(z - 1)^3(z + 2\mathrm{i})^2}\mathrm{d}z.
\]
\item[二]（20分）设\(U, V\)是复平面上区域，\(f : U \to V\)是连续函数，\(g : V \to
\mathbb{C}\)是非常值全纯函数，使得\(F = g \circ f : U \to \mathbb{C}\)也是全纯函数。
证明\(f\)是全纯函数。
\item[三]（20分）
将\(f(z) = (1 - z)^{-2}\)在\(z = 0\)处展开成幂级数，并求出该幂级数的收敛半径。
\item[四]（20分）求满足如下条件的\(\mathbb{C}\)上的半纯函数\(f\):
\[
    f\left( \frac{1}{n} \right) = \frac{n^3}{n-1},\quad n = 1, 2, \dots
\]
并求出\(f\)在\(\mathbb{C}\)中所有孤立奇点处的留数。
\item[五]（20分）设\(f\)是定义在\(0 < |z| < 1\)上的全纯函数，且存在常数\(C > 0\)和
正整数\(k\)使得
\[
    \int_{|z| = r}\left| z^k f(z) \right||\mathrm{d}z| \le C
\]
对所有的\(r \in (0, 1)\)成立。证明\(0\)是\(f\)的可去奇点或极点。
\end{enumerate}
